\documentclass[a4paper,oneside]{report}
% PT: notar que o relatório pode ser frente-e-verso ou apenas uma página por folha
% EN: note that the report can be double-sided or single-sided

\usepackage{sty/project-report} % formato ``oficial'' / official format
\reportlanguage{en} % Choose 'pt' or 'en'

\usepackage[english,portuguese]{babel} % Load both, last one is main
\ifthenelse{\equal{\selectedReportLanguage}{en}}
  {\AtBeginDocument{\selectlanguage{english}}}
  {\AtBeginDocument{\selectlanguage{portuguese}}}
\usepackage[T1]{fontenc}
\usepackage[utf8]{inputenc} % Unicode
\usepackage{graphicx} % incluir imagens / include graphics
\usepackage{url} % typeset URL's
\usepackage{textcomp} % símbolo do euro / euro symbol
\usepackage{multirow} % células com mais que uma linha em tabelas / multi-line cells in tables
\usepackage{booktabs}
\usepackage{tabularx}
\usepackage{amsmath}
\usepackage{tikz}
\usetikzlibrary{arrows.meta,positioning,calc,fit,shapes.geometric}
\usepackage{csquotes} % recommended by biblatex with babel
\usepackage[backend=biber,style=ieee]{biblatex}
\addbibresource{refs.bib}
\usepackage[pdftex,hidelinks]{hyperref}
\setlength{\emergencystretch}{3em} % fewer overfull lines with long identifiers
% Visible marker for items the author must complete before submission
\newcommand{\todo}[1]{\textbf{[TO DO: #1]}}
\usepackage[newfloat]{minted}
\setminted{breaklines,fontsize=\scriptsize,frame=single}

\begin{document}

\dept{School of Technology and Management (ESTiG)}
\course{Informatics Engineering}
\courseyear{2026}

\title{Cybersecurity Risks of Artificial Intelligence:\\[0.35em]
A Practical Study of Deepfakes,\\
AI-Generated Phishing, and Adversarial Attacks}
\author{Rinor Basholli}
\authnum{69492}

\supervisor{Tiago Pedrosa}

% Coloca a capa, primeira página e outros / Put the cover, first page and others

\beforepreface

\prefacesection*{Dedicatória/Dedication}
Dedico este relatório à minha família e ao Instituto Politécnico de Bragança.

I dedicate this report to my family and to the Instituto Politécnico de Bragança.

\prefacesection{Agradecimentos/Acknowledgments}
Gostaria de agradecer ao Professor Tiago Pedrosa pela orientação, pelas correcções directas e por exigir números que eu consiga defender. Agradeço também à ESTiG e ao Instituto Politécnico de Bragança pelas condições de trabalho.

I would like to thank Professor Tiago Pedrosa for the supervision, for the direct corrections, and for asking for numbers that I can defend. I also thank ESTiG and the Instituto Politécnico de Bragança for the working conditions.

\prefacesection{Resumo}
A inteligência artificial generativa reduz o custo do phishing, dos deepfakes e dos ataques adversariais, mas os estudantes encontram normalmente estes riscos em ferramentas separadas, que raramente distinguem resultados medidos de texto gerado por um modelo. Este relatório apresenta o InsightFlow AI, um laboratório web que reúne phishing gerado por IA, cenários de deepfake e ataques adversariais num único fluxo autenticado, e que identifica a proveniência de cada resultado analítico (\texttt{gemini}, \texttt{fallback}, \texttt{preset} ou \texttt{computed}). O laboratório é composto por um cliente React, por um servidor Express que invoca a API Gemini e aplica autenticação, limites de pedidos, limites de tamanho e mitigações de injeção de prompts, e por um serviço Python que executa os ataques FGSM e PGD sobre uma ResNet-56 pré-treinada no CIFAR-10. Foram realizadas duas avaliações reprodutíveis. Num corpus reconstruído de 187 emails institucionais, em que as mensagens legítimas e de phishing partilham a redação, um classificador TF-IDF com regressão logística baseado apenas no texto obteve um F1 de 0,708 em validação cruzada, abaixo de uma linha de base que classifica tudo como phishing (0,736); a inclusão do endereço do remetente elevou o F1 para 0,806 e a AUC para 0,837. Em 200 imagens de teste do CIFAR-10, a exatidão de 0,955 desceu para 0,05 com PGD a $\varepsilon = 2/255$ e para 0 a $8/255$, com todas as perturbações verificadas dentro do respetivo limite. Os 31 testes automáticos passam. Os resultados indicam que uma redação cuidada não é um sinal fiável de phishing e que classificadores sem defesas são frágeis; o efeito pedagógico do laboratório e a exatidão da análise de emails baseada em LLM ficam por avaliar.

\textbf{Palavras-chave:} cibersegurança, inteligência artificial, phishing, deepfakes, exemplos adversariais, educação em segurança.

\prefacesection{Abstract}
Generative artificial intelligence lowers the cost of phishing, deepfakes and adversarial attacks, yet students usually meet these risks in separate tools that rarely distinguish measured results from model-generated text. This report presents InsightFlow AI, a web laboratory that places AI-generated phishing, deepfake scenarios and adversarial attacks in one authenticated workflow and labels every analytical output with its provenance (\texttt{gemini}, \texttt{fallback}, \texttt{preset} or \texttt{computed}). The laboratory consists of a React client, an Express server that calls the Gemini API and enforces authentication, rate limits, request-size limits and prompt-injection mitigations, and a Python service that runs FGSM and PGD attacks on a pretrained CIFAR-10 ResNet-56. Two reproducible evaluations were carried out. On a reconstructed corpus of 187 institutional emails in which legitimate and phishing messages share their wording, a TF-IDF and logistic-regression classifier using the wording alone reached a cross-validated F1 of 0.708, below an always-phishing baseline of 0.736, whereas adding the sender address raised F1 to 0.806 and AUC to 0.837. On 200 CIFAR-10 test images, a clean accuracy of 0.955 fell to 0.05 under PGD at $\varepsilon = 2/255$ and to 0 at $8/255$, with every perturbation verified to lie within its budget. All 31 automated tests pass. The results indicate that polished wording is not a reliable phishing signal and that undefended classifiers are fragile; the educational effect of the laboratory and the accuracy of its LLM-based email analysis remain to be evaluated.

\textbf{Keywords:} cybersecurity, artificial intelligence, phishing, deepfakes, adversarial examples, security education.

% Coloca índices / Put indices
\afterpreface

\prefacesection{Lista de Abreviaturas/Abbreviations}
\begin{description}
 \item[AI] Artificial Intelligence
 \item[API] Application Programming Interface
 \item[ASR] Attack Success Rate
 \item[AUC] Area Under the (ROC) Curve
 \item[CIFAR] Canadian Institute for Advanced Research (image dataset)
 \item[CV] Cross-Validation
 \item[DKIM] DomainKeys Identified Mail
 \item[DMARC] Domain-based Message Authentication, Reporting, and Conformance
 \item[ESTiG] Escola Superior de Tecnologia e Gestão
 \item[FGSM] Fast Gradient Sign Method
 \item[FR] Functional Requirement
 \item[GDPR] General Data Protection Regulation
 \item[IPB] Instituto Politécnico de Bragança
 \item[JSON] JavaScript Object Notation
 \item[LLM] Large Language Model
 \item[NFR] Non-Functional Requirement
 \item[PGD] Projected Gradient Descent
 \item[ROC] Receiver Operating Characteristic
 \item[RQ] Research Question
 \item[SPF] Sender Policy Framework
 \item[STRIDE] Spoofing, Tampering, Repudiation, Information disclosure, Denial of service, Elevation of privilege
 \item[TF-IDF] Term Frequency--Inverse Document Frequency
\end{description}

\bodystart

% inclusão do texto, propriamente dito / include the text, properly speaking
\chapter{Introduction}
\label{cha:intro}

This chapter states the context of the work, the problem I tried to treat, the objectives, the research questions, the contributions, and the structure of the report. I wrote it as the author of the laboratory, not as an external auditor.

\section{Context}
\label{sec:context}

Artificial intelligence is now easy to call from a student project. A language model can write a fluent email in a few seconds~\cite{hazell2023spear,heiding2024devising}. A generative model can produce a face that never sat in front of a camera~\cite{karras2019stylegan,mirsky2021deepfake}. A small change in pixel values can move the class of an image classifier~\cite{szegedy2014intriguing,goodfellow2015explaining}. Those three facts are already in the public literature. What I did not have, as a student at ESTiG, was one place where I could show the three risks together and say honestly where each number came from.

The official title of this project is \emph{Cybersecurity Risks of Artificial Intelligence: A Practical Study of Deepfakes, AI-Generated Phishing, and Adversarial Attacks}. The live brand of the application is InsightFlow AI, also labelled AI Cybersecurity Risk Lab. I implemented a web laboratory for the Informatics Engineering project course at the Instituto Politécnico de Bragança, School of Technology and Management (ESTiG). Professor Tiago Pedrosa supervised the work.

I implemented the laboratory in this repository, added tests, and wired a local evaluation service. Appendix~\ref{app:ai} declares where a language model is used. I still have to defend every number in the viva.

\section{Problem}
\label{sec:problem}

Teaching material on phishing, deepfakes, and adversarial examples usually lives in separate tools. GoPhish is a campaign platform~\cite{gophish}. FaceForensics++ and Celeb-DF are research datasets~\cite{rossler2019faceforensics,li2020celebdf}. Torchattacks is a PyTorch library for FGSM and PGD~\cite{kim2020torchattacks,goodfellow2015explaining,madry2018towards}. I did not find a single educational web platform that covers the three topics in one student workflow and labels every analytical answer with a source field.

There is a second problem. It is easy to let a language model invent a confidence score. My supervisor asked that at least one module compute from real model outputs, and that I do not present dashboard decoration as science. The 75\% and 51\% figures on the dashboard come from score bands I wrote in \texttt{academicDataset.ts} with seed 20260927. They are configuration. They are not findings. Chapter~\ref{cha:eval} repeats that point.

\section{Objectives}
\label{sec:objectives}

I set the following objectives.

\begin{enumerate}
\item Build one educational web laboratory that shows AI-generated phishing, a deepfake teaching tab, and an adversarial image attack in the same application.
\item Keep a login lock on the laboratory in production (session cookie or HTTP Basic Auth), with secrets outside git.
\item Label every analytical JSON response with \texttt{source} equal to \texttt{gemini}, \texttt{fallback}, \texttt{preset}, or \texttt{computed}.
\item Use Gemini only on the teaching email routes (and on non-preset deepfake filenames). Do not use Gemini to invent adversarial confidences.
\item Implement FGSM and PGD against a pretrained classifier, with a path to export clean accuracy, robust accuracy, attack success rate, and perturbation size.
\item Write functional and non-functional requirements, use cases, architecture, tests, a STRIDE note, and an ethics note.
\item Cite only numbers that exist in exported files. If a file is missing, say so.
\end{enumerate}

\section{Research questions}
\label{sec:rqs}

I used three research questions. They are teaching questions, not claims about the internet at large.

\begin{description}
\item[RQ1] How can a student laboratory generate and analyse AI-written phishing mail next to conventional phishing and legitimate institutional text, while saying whether Gemini, a fallback, or a dataset label produced the answer?
\item[RQ2] How can a deepfake tab teach common visual and audio talking points without pretending to inspect media bytes?
\item[RQ3] How can FGSM and PGD be run on a real pretrained classifier so that clean accuracy and robust accuracy are meaningful, and so that Gemini is not used as a substitute when the attack service is down?
\end{description}

\section{Contributions}
\label{sec:contrib}

The main contribution is the laboratory itself: a React and Express application, a Python evaluation service, and this report.

I also contributed the following concrete pieces.

\begin{itemize}
\item An in-memory teaching set of 90 emails (30 legitimate, 30 conventional phishing, 30 AI-style phishing) generated with seed 20260927, plus three authored exemplars and four educational reconstructions of public scam or notice patterns. I did not harvest a mailbox.
\item Express routes that return \texttt{source} on every analysis. The adversarial route returns 503 if the Python service is down. It does not call Gemini.
\item HTTP Basic Auth and a login cookie, a 32\,kB JSON body limit, a prompt-injection reject list, and an hourly rate limit (10 requests per IP for \texttt{/api/generate-phishing}, 40 for other \texttt{POST /api/} paths).
\item Automated tests in \texttt{tests/labSecurity.test.ts} and \texttt{tests/apiContract.test.ts}.
\item A phishing benchmark script that exported \texttt{eval/phishing/results/phishing\_metrics.json}. I report those numbers in Chapter~\ref{cha:eval}. The first F1 $= 1.0$ run was a construction error. The cited file reports wording-only and with-sender side by side, plus a dummy baseline and 5-fold CV.
\item An adversarial batch that exported \texttt{eval/adversarial/results/adversarial\_metrics.json} for $n=200$, seed 20260927, CIFAR-10 ResNet-56. Chapter~\ref{cha:eval} copies that file.
\end{itemize}

\section{Structure of the report}
\label{sec:structure}

Chapter~\ref{cha:sota} reviews related tools and papers and states the gap I tried to fill. Chapter~\ref{cha:req} lists requirements and use cases. Chapter~\ref{cha:arch} describes the architecture, sequence flows, the client-only data model, and the API contracts. Chapter~\ref{cha:impl} explains the implementation choices and the tests. Chapter~\ref{cha:eval} describes the evaluation method and the numbers I actually have. Chapter~\ref{cha:disc} discusses limitations, ethics, dual use, and STRIDE. Chapter~\ref{cha:conc} concludes and lists future work. Appendix~\ref{app:ai} is the AI declaration. Appendix~\ref{app:eval} says how to run the evaluation scripts.

\chapter{State of the art}
\label{cha:sota}

This chapter reviews the work on which the laboratory builds. It is organised by theme: AI-generated phishing and its detection (Section~\ref{sec:sota-phish}), deepfakes (Section~\ref{sec:sota-df}), adversarial examples (Section~\ref{sec:sota-adv}), the security of applications that embed LLMs (Section~\ref{sec:sota-llmsec}) and security education (Section~\ref{sec:sota-edu}). Section~\ref{sec:gap} compares existing tools and states the gap that this work addresses.

\section{AI-generated phishing and its detection}
\label{sec:sota-phish}

\subsection{Phishing as an attack vector}
Phishing remains one of the most common initial-access techniques in the threat landscape described by the European Union Agency for Cybersecurity~\cite{enisa2024threat}, and the Anti-Phishing Working Group continues to record high volumes of phishing attacks~\cite{apwg2024}. Conventional lures rely on recognisable cues: artificial urgency, generic greetings, grammatical errors, look-alike domains and requests for credentials or payment.

\subsection{Phishing generated by language models}
LLMs remove several of those cues. Hazell showed that commercially available LLMs can generate personalised spear-phishing messages for hundreds of targets at negligible cost~\cite{hazell2023spear}. Heiding et al.\ compared LLM-generated phishing emails with emails crafted by human experts in a user study, and also evaluated LLMs as detectors of malicious intent in emails~\cite{heiding2024devising}. The common conclusion is that fluent, well-tailored wording can no longer be taken as evidence that a message is legitimate.

\subsection{Machine-learning detection}
Automatic phishing detection predates LLMs. Fette et al.\ classified emails with a random forest over structural features such as IP-based URLs, link text that does not match its target, domain age and the output of a spam filter, and reported detection of more than 96\% of phishing emails with a false-positive rate of about 0.13\% on their dataset~\cite{fette2007learning}. A later systematic review by Salloum et al.\ surveys natural-language-processing approaches, ranging from TF-IDF features with classical classifiers to deep neural models~\cite{salloum2022phishing}. A TF-IDF representation with logistic regression is a standard, interpretable baseline in that body of work; it is the baseline adopted in Chapter~\ref{cha:eval}.

LLMs have also been proposed as detectors. Koide et al.\ converted emails into prompts and reported 99.70\% accuracy with GPT-4 on a dataset of 2{,}010 emails, ahead of other LLMs and conventional baselines~\cite{koide2024chatspamdetector}. Such results depend strongly on the dataset; they show that an LLM can produce a verdict and a rationale, which is the design of the email-analysis module in this work, but they do not transfer to other corpora without re-evaluation.

\subsection{Sender authentication}
If wording is no longer a reliable signal, the authenticity of the sender becomes central. Three IETF standards address it: Sender Policy Framework (SPF) authorises sending hosts for a domain~\cite{rfc7208}; DomainKeys Identified Mail (DKIM) signs messages with a domain key~\cite{rfc6376}; and Domain-based Message Authentication, Reporting, and Conformance (DMARC) lets a domain owner publish how receivers should treat mail that fails these checks~\cite{rfc9989}. These mechanisms verify the sending domain, not the content, and are therefore unaffected by how convincing an LLM makes the text.

\section{Deepfakes}
\label{sec:sota-df}

Mirsky and Lee survey the creation and detection of deepfakes of the human face and body, including reenactment and replacement techniques~\cite{mirsky2021deepfake}. Generative adversarial networks such as StyleGAN produce high-resolution synthetic faces~\cite{karras2019stylegan}. Detection research relies on benchmark datasets: FaceForensics++ provides manipulated videos produced with several methods~\cite{rossler2019faceforensics}, Celeb-DF provides higher-quality forgeries that are harder to detect~\cite{li2020celebdf}, and the DeepFake Detection Challenge released a large dataset with a public competition~\cite{dolhansky2020dfdc}.

Early detectors exploited specific artefacts. Li, Chang and Lyu observed that early face-swap videos showed abnormal eye blinking, because training images rarely contain closed eyes, and detected fakes from blink patterns~\cite{li2018ictu}. Such cues age quickly as generators improve. For synthetic speech, Müller et al.\ re-evaluated audio deepfake detectors under uniform conditions and found that their performance degrades sharply on real-world recordings outside the ASVspoof benchmark~\cite{muller2022audio}. This limited generalisation is one reason why this work does not ship a detector (Section~\ref{sec:impl-df}).

A complementary approach is provenance: the Coalition for Content Provenance and Authenticity specifies cryptographically signed metadata, known as Content Credentials, that records how a media file was created and edited~\cite{c2pa2spec}. Provenance does not detect manipulation; it allows authentic media to be verified.

\section{Adversarial examples}
\label{sec:sota-adv}

\subsection{Attacks}
Szegedy et al.\ showed that deep networks misclassify inputs that differ from correctly classified ones by small, carefully chosen perturbations~\cite{szegedy2014intriguing}. Goodfellow, Shlens and Szegedy proposed the Fast Gradient Sign Method (FGSM), a single step of size $\varepsilon$ in the direction of the sign of the input gradient of the loss~\cite{goodfellow2015explaining}:
\begin{equation}
x_{\mathrm{adv}} = \mathrm{clip}_{[0,1]}\bigl(x + \varepsilon \cdot \mathrm{sign}(\nabla_x J(\theta, x, y))\bigr).
\label{eq:fgsm}
\end{equation}
Madry et al.\ formulated robustness as a min--max problem and used Projected Gradient Descent (PGD), an iterated FGSM with projection back onto the $\varepsilon$-ball, as a strong first-order attack~\cite{madry2018towards}. Biggio and Roli review a decade of adversarial machine learning, including evasion and poisoning attacks~\cite{biggio2018wild}. The Torchattacks library implements these attacks for PyTorch~\cite{kim2020torchattacks}.

\subsection{Evaluation and defences}
Evaluating robustness is error-prone. Carlini et al.\ set out principles for sound evaluations, including a precise threat model, strong adaptive attacks and sanity checks such as accuracy approaching zero for large budgets~\cite{carlini2019evaluating}. Input transformations such as JPEG compression, bit-depth reduction and total-variance minimisation have been proposed as defences~\cite{guo2018countering}; however, Athalye, Carlini and Wagner showed that defences relying on obfuscated gradients can be circumvented by adapted attacks, defeating six of seven such defences completely~\cite{athalye2018obfuscated}. The adversarial module of this work therefore presents preprocessing only as a defence to be tested, not as a guarantee.

\subsection{Model choice}
Residual networks~\cite{he2016resnet} are a standard architecture in this literature. CIFAR-10 consists of $32 \times 32$ colour images in ten classes~\cite{krizhevsky2009cifar}. Pretrained CIFAR-10 residual networks are distributed in the \texttt{pytorch-cifar-models} repository~\cite{chenyaofo2021cifar}, which provides the model used in this work.

\section{Security of LLM-integrated applications}
\label{sec:sota-llmsec}

An application that inserts user text into an LLM prompt inherits the risk of prompt injection. Perez and Ribeiro showed that simple handcrafted inputs can hijack a model's goal or leak its prompt~\cite{perez2022ignore}. Greshake et al.\ introduced indirect prompt injection, in which instructions are hidden in data that the application later passes to the model rather than in the user's own request, and demonstrated it against deployed systems~\cite{greshake2023signed}. An email-analysis service is an instance of the indirect case, because the email body is chosen by a third party. These works motivate the input handling described in Section~\ref{sec:limits}.

\section{Security education}
\label{sec:sota-edu}

Research on anti-phishing education indicates that interactive training can outperform passive material. Sheng et al.\ evaluated Anti-Phishing Phil, an online game, in a 42-participant laboratory study and found that participants who played the game identified phishing websites better than those who read existing online training material~\cite{sheng2007phil}. Kumaraguru et al.\ studied anti-phishing training designs and their evaluation in more depth~\cite{kumaraguru2010johnny}. Practitioner platforms such as TryHackMe and Hack The Box offer hands-on security exercises across many topics~\cite{tryhackme,hackthebox}. These platforms are broad and hosted by third parties; they are not small applications that a student can deploy, inspect and explain in full.

\section{Comparison and gap}
\label{sec:gap}

Table~\ref{tab:gap} compares representative tools and resources against three criteria derived from the problem statement in Section~\ref{sec:problem}: coverage of the three risks, availability as a self-hostable teaching application, and transparency about whether an output is measured, generated or configured.

\begin{table}[htbp]
\centering
\caption{Comparison of related tools and resources.}
\label{tab:gap}
\scriptsize
\begin{tabular}{p{2.6cm}p{1.7cm}p{1.7cm}p{1.7cm}p{1.5cm}p{1.8cm}}
\toprule
Tool or resource & Phishing & Deepfakes & Adversarial & Self-hostable teaching app & Output provenance shown \\
\midrule
GoPhish~\cite{gophish} & Campaigns & --- & --- & Yes & Not applicable \\
FaceForensics++, Celeb-DF~\cite{rossler2019faceforensics,li2020celebdf} & --- & Datasets & --- & No & Ground-truth labels \\
Torchattacks~\cite{kim2020torchattacks} & --- & --- & Library & No (library) & Computed \\
Anti-Phishing Phil~\cite{sheng2007phil} & URL training & --- & --- & Game & Not applicable \\
TryHackMe, Hack The Box~\cite{tryhackme,hackthebox} & Some rooms & Rare & Rare & No (hosted) & Varies \\
ChatSpamDetector~\cite{koide2024chatspamdetector} & LLM detection & --- & --- & Research tool & LLM rationale \\
\textbf{InsightFlow AI (this work)} & Generate and analyse & Scenarios (no media) & FGSM/PGD & Yes & Per response \\
\bottomrule
\end{tabular}
\end{table}

Each risk is well covered on its own: phishing training and detection, deepfake datasets and detectors, and adversarial-attack libraries are all mature. No resource identified in the search combines the three risks in one self-hostable application for teaching, and the general-purpose platforms do not distinguish, per response, between a computed result, an LLM's narration and a value configured by the tool's author. The gap addressed by this work is therefore one of integration and transparency for teaching, not a new attack or detection technique.

\section{Summary}
\label{sec:sota-summary}

LLMs make phishing wording an unreliable signal, which shifts attention to sender authentication; deepfake detectors generalise poorly, which favours provenance and awareness; and adversarial robustness must be evaluated with strong attacks and verified budgets. Chapter~\ref{cha:req} turns these observations into requirements for the laboratory.

\chapter{Requirements and use cases}
\label{cha:req}

This chapter specifies what the laboratory must do. It identifies the actors, lists the functional and non-functional requirements with their origin and priority, and details the main use cases. The requirements were refined iteratively during development; the \emph{Source} column records why each one exists, so that it can be traced to the problem statement, to the threat analysis or to the evaluation design.

\section{Actors}
\label{sec:actors}

Three human actors and two external systems interact with the laboratory (Figure~\ref{fig:usecases}).

\begin{description}
\item[Learner] A student who uses the laboratory modules to generate and inspect phishing emails, examine deepfake scenarios and run adversarial attacks.
\item[Instructor] A teacher or examiner who demonstrates the modules, inspects the provenance of results and consults the exported evaluation files.
\item[Operator] The person who deploys the laboratory, sets the access credentials and API keys, and starts the attack service.
\item[Gemini API] The external LLM service used for phishing generation, email analysis and deepfake narration.
\item[Attack service] The local Python service that runs FGSM and PGD on the pretrained classifier.
\end{description}

\begin{figure}[htbp]
\centering
\begin{tikzpicture}[
  actor/.style={draw, rounded corners, font=\scriptsize\bfseries, minimum width=2.0cm, minimum height=0.7cm, align=center},
  uc/.style={draw, ellipse, font=\scriptsize, minimum width=3.6cm, minimum height=0.75cm, align=center},
  line/.style={thin}
]
\draw[rounded corners] (2.3,1.0) rectangle (8.9,-7.6);
\node[font=\scriptsize\bfseries] at (5.6,0.7) {InsightFlow AI};
\node[uc] (u1) at (5.6,0.0) {UC1 Sign in};
\node[uc] (u2) at (5.6,-1.1) {UC2 Load teaching set};
\node[uc] (u3) at (5.6,-2.2) {UC3 Generate phishing email};
\node[uc] (u4) at (5.6,-3.3) {UC4 Analyse email};
\node[uc] (u5) at (5.6,-4.4) {UC5 Inspect deepfake scenario};
\node[uc] (u6) at (5.6,-5.5) {UC6 Run adversarial attack};
\node[uc] (u7) at (5.6,-6.6) {UC7 Reproduce evaluation};
\node[actor] (l) at (0,-1.6) {Learner};
\node[actor] (i) at (0,-4.6) {Instructor};
\node[actor] (o) at (0,-6.6) {Operator};
\node[actor] (g) at (11.2,-2.7) {Gemini API};
\node[actor] (a) at (11.2,-5.5) {Attack service};
\foreach \u in {u1,u2,u3,u4,u5,u6} \draw[line] (l) -- (\u.west);
\foreach \u in {u1,u5,u6,u7} \draw[line] (i) -- (\u.west);
\draw[line] (o) -- (u7.west);
\draw[line] (u3.east) -- (g); \draw[line] (u4.east) -- (g); \draw[line] (u5.east) -- (g);
\draw[line] (u6.east) -- (a);
\end{tikzpicture}
\caption{Use-case diagram of the laboratory.}
\label{fig:usecases}
\end{figure}

\section{Functional requirements}
\label{sec:fr}

Table~\ref{tab:fr} lists the functional requirements, prioritised with the MoSCoW scheme (\emph{Must}, \emph{Should}, \emph{Could}). The source column refers to the problem statement (P), the threat analysis in Section~\ref{sec:stride} (T) or the evaluation design (E).

\begin{table}[htbp]
\centering
\caption{Functional requirements.}
\label{tab:fr}
\scriptsize
\begin{tabular}{P{0.7cm}P{6.4cm}P{1.0cm}P{0.7cm}P{2.7cm}}
\toprule
ID & Requirement & Priority & Source & Realised by \\
\midrule
FR1 & Require authentication for every API call when credentials are configured, and refuse to serve the laboratory in production when they are not. & Must & T & \texttt{labSecurity.ts} \\
FR2 & Generate a teaching phishing email from a scenario, a topic, optional instructions and a complexity level. & Must & P & \path{/api/generate-phishing} \\
FR3 & Analyse a subject, a body and an optional sender, and return scores, indicators, mitigations and a predicted class with a rationale. & Must & P & \path{/api/analyze-email} \\
FR4 & Accept a local \texttt{.txt} or \texttt{.eml} file and extract sender, recipient, subject and body in the browser. & Should & P & \texttt{parseEmailFile} \\
FR5 & Load a seeded, reproducible set of 90 teaching emails (30 per class) into the dashboard. & Must & P & \texttt{generate90Samples} \\
FR6 & Analyse a deepfake \emph{scenario} from a filename, a media type and a description, without reading media content. & Must & P & \path{/api/analyze-media-deepfake} \\
FR7 & Return fixed teaching content when the filename matches one of three presets. & Should & P & \texttt{PRE\_BAKED\_DEEPFAKES} \\
FR8 & Run FGSM or PGD on a pretrained CIFAR-10 classifier for a chosen class and budget, untargeted from the user interface and optionally targeted through the API. & Must & P, E & \path{/api/adversarial-perturb}, \texttt{service.py} \\
FR9 & Label every successful analytical response with \texttt{source} $\in$ \{\texttt{gemini}, \texttt{fallback}, \texttt{preset}, \texttt{computed}\}. & Must & P & \texttt{server.ts}, \texttt{service.py} \\
FR10 & Return a deterministic fallback when Gemini is unavailable on the email and deepfake routes; return HTTP~503 on the adversarial route when the attack service is unavailable. & Must & P & Route handlers \\
FR11 & Expose the health of the attack service and the exported metric files. & Should & E & \path{/api/eval/*} \\
FR12 & Present dashboard aggregates as configured teaching values, not as measurements. & Must & P & \path{UnifiedAnalysisDashboard.tsx} \\
FR13 & Provide scripts that regenerate the phishing benchmark and the adversarial batch results from code. & Must & E & \path{eval/} \\
\bottomrule
\end{tabular}
\end{table}

\section{Non-functional requirements}
\label{sec:nfr}

Table~\ref{tab:nfr} lists the non-functional requirements. Where possible, each has a measurable target that is verified in Chapter~\ref{cha:eval}.

\begin{table}[htbp]
\centering
\caption{Non-functional requirements with measurable targets.}
\label{tab:nfr}
\scriptsize
\begin{tabular}{P{0.9cm}P{2.1cm}P{6.6cm}P{0.8cm}}
\toprule
ID & Quality & Requirement and target & Source \\
\midrule
NFR1 & Security & JSON request bodies larger than 32\,kB are rejected. & T \\
NFR2 & Security & At most 10 \texttt{POST} requests per hour per client address to the generation route, and 40 to every other \texttt{POST /api/} route. & T \\
NFR3 & Security & Credentials are compared in constant time; empty or length-mismatched values never authenticate. & T \\
NFR4 & Security & Sessions expire after 12 hours; session tokens carry 192 bits of randomness and are sent in \texttt{HttpOnly}, \texttt{SameSite=Lax} cookies. & T \\
NFR5 & Security & Secrets are supplied through environment variables and are never committed to the repository. & T \\
NFR6 & Integrity of results & Every adversarial example satisfies $\|x_{\mathrm{adv}} - x\|_\infty \le \varepsilon$; a violation is reported as an error, never returned as a result. & E \\
NFR7 & Reproducibility & Running each evaluation script twice on the same commit yields identical metric files (fixed seed 20260927). & E \\
NFR8 & Availability & The website remains usable when Gemini or the attack service is down; affected routes degrade as specified in FR10. & P \\
NFR9 & Testability & The automated test suite runs offline, without calling Gemini, in under 10 seconds. & E \\
NFR10 & Portability & The application can be built and run as a container image. & P \\
NFR11 & Security & At most 10 failed sign-in attempts per client address within 15 minutes; further attempts receive HTTP~429. & T \\
\bottomrule
\end{tabular}
\end{table}

\section{Main use cases}
\label{sec:uc}

Tables~\ref{tab:uc1}--\ref{tab:uc6} detail the three use cases that carry the main security and evaluation logic. The remaining use cases follow the same pattern and are covered by the sequence diagrams in Chapter~\ref{cha:arch}.

\begin{table}[htbp]
\centering
\caption{UC1 --- Sign in.}
\label{tab:uc1}
\scriptsize
\begin{tabular}{P{2.4cm}P{10.0cm}}
\toprule
Actor & Learner or Instructor \\
Precondition & The Operator has set \texttt{LAB\_BASIC\_USER} and \texttt{LAB\_BASIC\_PASS}. \\
Main flow & (1) The browser loads the login page, which is served without a session. (2) The actor submits a username and password. (3) The server compares both values in constant time. (4) The server creates a session and returns an \texttt{HttpOnly} cookie. (5) The application shows the dashboard. \\
Alternative flows & (3a) Wrong credentials: HTTP~401, the form shows an error; after 10 failures from the same address within 15 minutes, HTTP~429. (1a) Credentials not configured in production: every request receives HTTP~503 and the laboratory stays closed. (1b) Non-interactive clients may use HTTP Basic authentication instead of the cookie. \\
Postcondition & Subsequent \path{/api/*} calls are accepted until the session expires or the actor signs out. \\
\bottomrule
\end{tabular}
\end{table}

\begin{table}[htbp]
\centering
\caption{UC4 --- Analyse email.}
\label{tab:uc4}
\scriptsize
\begin{tabular}{P{2.4cm}P{10.0cm}}
\toprule
Actor & Learner; Gemini API \\
Precondition & The actor is signed in. \\
Main flow & (1) The actor pastes a subject, body and optional sender, or loads a \texttt{.txt}/\texttt{.eml} file that is parsed in the browser. (2) The browser posts the fields to \path{/api/analyze-email}. (3) The server places the email between explicit untrusted-data markers in the prompt and calls Gemini with a JSON response schema. (4) The server normalises the predicted class and returns the analysis with \texttt{source=gemini}. (5) The dashboard adds the result to the session history. \\
Alternative flows & (3a) Gemini fails: the server returns a keyword-heuristic analysis with \texttt{source=fallback}. (2a) More than 40 requests in the last hour: HTTP~429. (2b) Body larger than 32\,kB: HTTP~413. \\
Postcondition & The result is visible with its provenance; nothing is stored on the server. \\
\bottomrule
\end{tabular}
\end{table}

\begin{table}[htbp]
\centering
\caption{UC6 --- Run adversarial attack.}
\label{tab:uc6}
\scriptsize
\begin{tabular}{P{2.4cm}P{10.0cm}}
\toprule
Actor & Learner or Instructor; Attack service \\
Precondition & The actor is signed in; the Operator has started the attack service. \\
Main flow & (1) The actor selects a CIFAR-10 class, FGSM or PGD, and a budget $\varepsilon \in (0, 0.25]$. (2) The server forwards the request to the attack service. (3) The service selects the first test image of that class that the model classifies correctly, runs an untargeted attack, and measures $\|x_{\mathrm{adv}} - x\|_\infty$. (4) The service returns the clean and adversarial predictions with their softmax confidences, the budget and the measured perturbation with \texttt{source=computed}. \\
Alternative flows & (2a) Attack service unreachable: HTTP~503 with \texttt{source=unavailable}; no values are shown. (3a) Unknown class or budget outside the allowed range: the service answers HTTP~400 and the server reports HTTP~502. (3b) Measured perturbation exceeds $\varepsilon$: the service answers HTTP~500 and the server reports HTTP~502; no result is returned. (1a) Through the API only, a request with \texttt{targeted: true} runs a targeted attack towards a second class and reports whether the target was reached. \\
Postcondition & The displayed numbers are model outputs for one image; they are distinct from the batch evaluation in Chapter~\ref{cha:eval}. \\
\bottomrule
\end{tabular}
\end{table}

\section{Out of scope}
\label{sec:oos}

The following are deliberately excluded: sending email, harvesting or storing real mailboxes, detecting manipulation in uploaded media, training or fine-tuning models, per-user accounts and roles, and server-side storage of generated or analysed content. Chapter~\ref{cha:disc} discusses the consequences of these exclusions.

\section{Summary}

The requirements combine teaching functions (FR2--FR8), a provenance contract (FR9, FR10, FR12), security controls derived from the threat analysis (FR1, NFR1--NFR5, NFR11) and evaluation properties (FR13, NFR6, NFR7). Chapter~\ref{cha:arch} shows how the architecture realises them.

\chapter{Architecture}
\label{cha:arch}

This chapter describes how I organised the system after I had a working React laboratory and added a real evaluation path. I justified each split by what I needed to demonstrate, not by a catalogue of patterns.

\section{Overall architecture}
\label{sec:arch-overview}

Figure~\ref{fig:arch} shows the four processes I actually run.

\begin{figure}[htbp]
\centering
\begin{tikzpicture}[
  box/.style={draw, rounded corners, align=center, inner sep=6pt, font=\scriptsize, minimum width=3.1cm, minimum height=1.15cm},
  arr/.style={-{Stealth}, thick}
]
\node[box] (ui) at (0,0) {React UI\\Vite + four tabs};
\node[box] (ex) at (5.1,0) {Express\\ \texttt{server.ts}};
\node[box] (gem) at (5.1,2.3) {Gemini API\\ \texttt{gemini-3.5-flash}};
\node[box] (py) at (10.2,0) {Python eval\\FastAPI :8001};
\node[box] (m) at (10.2,-2.3) {CIFAR-10\\ResNet-56 weights};
\draw[arr] (ui) -- node[above, font=\scriptsize] {cookie / Basic + JSON} (ex);
\draw[arr] (ex) -- node[right, font=\scriptsize] {email / deepfake} (gem);
\draw[arr] (ex) -- node[above, font=\scriptsize] {\texttt{/attack}, \texttt{/health}} (py);
\draw[arr] (py) -- (m);
\end{tikzpicture}
\caption{Processes I used: the browser talks only to Express. Express talks to Gemini for teaching routes and to the Python service for real attacks.}
\label{fig:arch}
\end{figure}

I chose this split for a simple reason. The first application already had a React user interface and Express routes. Rewriting it as a notebook would have removed the laboratory that the examiners will click. Putting PyTorch inside Node would have been fragile. A second process on port 8001 lets me start or stop the attack service without taking the website down.

I run the Node application locally with \texttt{npm run dev} at \texttt{http://localhost:3000}. The real attack needs the Python service on port 8001.

\section{Client screens}
\label{sec:screens}

\texttt{App.tsx} keeps one \texttt{activeScreen} value: dashboard, phishing, deepfake, or adversarial. I chose a sidebar instead of separate deployments because the point of the project is one workflow.

The dashboard reads \texttt{phishingHistory}. It does not call the Python service. The phishing lab calls \texttt{/api/generate-phishing} and \texttt{/api/analyze-email}. The deepfake lab calls \texttt{/api/analyze-media-deepfake}. The attack lab calls \texttt{/api/adversarial-perturb}.

\section{Sequence: generate phishing}
\label{sec:seq-phish}

Figure~\ref{fig:seq-phish} is the generate path.

\begin{figure}[htbp]
\centering
\begin{tikzpicture}[
  every node/.style={font=\scriptsize},
  obj/.style={font=\scriptsize\bfseries},
  arr/.style={-{Stealth}}
]
\node[obj] (c) at (0,0) {Browser};
\node[obj] (s) at (4.0,0) {Express};
\node[obj] (g) at (8.0,0) {Gemini};
\draw (0,-0.3) -- (0,-5.4);
\draw (4.0,-0.3) -- (4.0,-5.4);
\draw (8.0,-0.3) -- (8.0,-5.4);
\draw[arr] (0,-0.8) -- node[above] {POST /api/generate-phishing} (4.0,-0.8);
\draw[arr] (4.0,-1.5) -- node[above] {injection / rate limit} (0,-1.5);
\node at (2.0,-1.9) {400 or 429 if rejected};
\draw[arr] (4.0,-2.5) -- node[above] {generateContent} (8.0,-2.5);
\draw[arr] (8.0,-3.3) -- node[above] {JSON email} (4.0,-3.3);
\draw[arr] (4.0,-4.1) -- node[above] {source=gemini} (0,-4.1);
\draw[arr] (4.0,-4.9) -- node[above] {on error: source=fallback} (0,-4.9);
\end{tikzpicture}
\caption{Generate-phishing sequence. A rejected prompt never reaches Gemini. A Gemini failure still returns a template.}
\label{fig:seq-phish}
\end{figure}

I kept the fallback because a viva demonstration should not die on a 503 from the model API. The source field is what stops that fallback from looking like a model measurement.

The analyse-email sequence is the same shape. The request body is \texttt{subject}, \texttt{body}, and optional \texttt{sender}. The success payload is \texttt{analysis} plus \texttt{source}.

\section{Sequence: deepfake scenario}
\label{sec:seq-df}

I implemented a branch before Gemini.

\begin{enumerate}
\item The browser posts \texttt{fileName}, \texttt{mediaType}, and \texttt{promptAnalysis}.
\item Express lowercases the filename. If it equals \texttt{ceo\_keynote\_urgent\_wire.mp4}, \texttt{financial\_analyst\_voicemail.wav}, or \texttt{id\_verification\_passport.png}, it returns the matching object from \texttt{PRE\_BAKED\_DEEPFAKES} with \texttt{source: "preset"}.
\item Otherwise it calls Gemini. On success the source is \texttt{gemini}. On error the source is \texttt{fallback} and \texttt{getDynamicFallbackDeepfake} builds a story from the filename and the word ``urgent''.
\end{enumerate}

No step reads an uploaded byte stream. I chose filename matching because it is deterministic for the three viva presets and because I did not have a licensed detector I could defend.

\section{Sequence: adversarial attack}
\label{sec:seq-adv}

Figure~\ref{fig:seq-adv} is the only path that I treat as computed.

\begin{figure}[htbp]
\centering
\begin{tikzpicture}[
  every node/.style={font=\scriptsize},
  obj/.style={font=\scriptsize\bfseries},
  arr/.style={-{Stealth}}
]
\node[obj] (c) at (0,0) {Browser};
\node[obj] (s) at (3.7,0) {Express};
\node[obj] (p) at (7.6,0) {FastAPI};
\node[obj] (m) at (11.0,0) {ResNet-56};
\draw (0,-0.3) -- (0,-5.0);
\draw (3.7,-0.3) -- (3.7,-5.0);
\draw (7.6,-0.3) -- (7.6,-5.0);
\draw (11.0,-0.3) -- (11.0,-5.0);
\draw[arr] (0,-0.9) -- node[above] {POST /api/adversarial-perturb} (3.7,-0.9);
\draw[arr] (3.7,-1.8) -- node[above] {POST /attack + token} (7.6,-1.8);
\draw[arr] (7.6,-2.7) -- node[above] {forward + FGSM/PGD} (11.0,-2.7);
\draw[arr] (11.0,-3.5) -- node[above] {softmax} (7.6,-3.5);
\draw[arr] (7.6,-4.3) -- node[above] {source=computed or 503} (0,-4.3);
\end{tikzpicture}
\caption{Adversarial sequence. If FastAPI is down, Express returns 503. Gemini is not called.}
\label{fig:seq-adv}
\end{figure}

\texttt{getAdversarialFallback} still exists in \texttt{server.ts}. I did not wire it to the POST handler. I left it unused on purpose so that a missing eval service cannot look like a successful attack.

\section{Data model}
\label{sec:data-model}

There is no database. I stored laboratory state in React \texttt{useState} hooks. Figure~\ref{fig:data} lists the fields.

\begin{figure}[htbp]
\centering
\begin{tikzpicture}[
  card/.style={draw, rectangle, align=left, font=\scriptsize, inner sep=6pt}
]
\node[card] (a) {
\textbf{App state (useState only)}\\[2pt]
activeScreen: dashboard $|$ phishing $|$ deepfake $|$ adversarial\\
lastPhishingAnalysis: object or null\\
phishingHistory: AnalyzedSample[\,]\\[4pt]
\textbf{AnalyzedSample}\\[2pt]
id, sender, recipient, subject, body\\
category: legitimate $|$ conventional\_phishing $|$ ai\_phishing\\
tag, realismScore, detectabilityScore, languagePolish\\
urgencyLevel, detectionDifficulty\\
indicators[], mitigationRules[], detailedAnalysis\\
analysisSource?: dataset $|$ gemini $|$ fallback\\
labelOrigin?: dataset $|$ known-sample $|$ llm $|$ fallback
};
\end{tikzpicture}
\caption{Client data model. Clearing the dataset only empties the React array.}
\label{fig:data}
\end{figure}

I chose client state because the 90 rows are generated in the browser from a seed. A database would have implied persistence that I do not have. The footer says the session is in memory. There is no PostgreSQL or SQLite.

The Python side keeps CIFAR-10 on disk under \texttt{eval/adversarial/results/cifar10/} after the first download. That cache is local data for the model, not a user database.

\section{API contracts}
\label{sec:api}

All JSON bodies are limited to 32\,kB. POST routes under \texttt{/api/} are rate limited. Table~\ref{tab:api} is the contract I implemented.

\begin{table}[htbp]
\centering
\caption{HTTP contracts for the four POST routes and the evaluation health route.}
\label{tab:api}
\scriptsize
\begin{tabular}{p{3.4cm}p{4.2cm}p{4.6cm}}
\toprule
Route & Request body & Success / error \\
\midrule
\texttt{POST /api/generate-phishing} & \texttt{scenario}, \texttt{promptText}, \texttt{customPrompt}, \texttt{complexity} & \texttt{\{success, source, email\}} with subject, body, realismScore, spoofingIndicators, mitigationRules. 400 if injected. 429 if limited. \\
\texttt{POST /api/analyze-email} & \texttt{subject}, \texttt{body}, \texttt{sender} & \texttt{\{success, source, analysis\}} with scores, predictedCategory, categoryRationale, indicators, mitigationRules, detailedAnalysis. \\
\texttt{POST /api/analyze-media-deepfake} & \texttt{fileName}, \texttt{mediaType}, \texttt{promptAnalysis} & \texttt{\{success, source, deepfakeAnalysis\}} with facialIrregularities, audioAnomalies, simulatedClassification, markers, actionableDefences. \\
\texttt{POST /api/adversarial-perturb} & \texttt{originalClass}, \texttt{targetClass}, \texttt{noiseLevel}, \texttt{architecture} & \texttt{\{success, source: computed, perturbation\}} or 502/503. No Gemini fallback. \\
\texttt{GET /api/eval/health} & none & Forwards Python \texttt{/health}. 503 if the service is down. \\
\bottomrule
\end{tabular}
\end{table}

I also added two GET helpers that only send a file if it exists: \texttt{/api/eval/phishing} and \texttt{/api/eval/adversarial}. They are not extra attack surfaces in the teaching sense. They just expose the JSON I need for Chapter~\ref{cha:eval}.

The Python \texttt{POST /attack} body is the same class pair plus \texttt{method} (default \texttt{fgsm}). If \texttt{EVAL\_INTERNAL\_TOKEN} is set, Express sends it as \texttt{X-Eval-Token}.

\section{Security controls in the request path}
\label{sec:arch-sec}

I placed three Express middleware checks in front of the routes: JSON size, \texttt{requireLabAccess}, and \texttt{rateLimitApi}. The generate-phishing handler then calls \texttt{promptLooksInjected} on the custom prompt and the prompt text. Chapter~\ref{cha:disc} maps those controls to STRIDE.

\chapter{Implementation}
\label{cha:impl}

This chapter explains how I built the laboratory and why I made each choice. I write in the past tense because this is the version I am submitting, not a plan.

\section{Starting point}
\label{sec:start}

I implemented the laboratory in React, Vite, TypeScript and Express: four tabs, authentication, source labels, tests, and a Python attack path. Appendix~\ref{app:ai} states where Gemini is used.

\section{Runtime and model identifiers}
\label{sec:runtime}

The Node process loads \texttt{.env} and then \texttt{.env.local}. The model id is:

\begin{minted}{javascript}
const GEMINI_MODEL = process.env.GEMINI_MODEL || 'gemini-3.5-flash';
\end{minted}

I used \texttt{gemini-3.5-flash} as the default because that is the id I configured for the teaching calls~\cite{geminiapi}. I can override it with \texttt{GEMINI\_MODEL} without editing \texttt{server.ts}. The process logs the id at startup. I still have to confirm, in the viva, which id the live key actually served on a given day. One labelled Gemini session is stored under \texttt{eval/gemini\_runs/}: three teaching routes with \texttt{source=gemini}. Live FGSM/PGD calls are stored under \texttt{eval/adversarial/results/live\_*.json}. I do not invent extra fallback counts beyond those files.

I chose the official \texttt{@google/genai} client and JSON response schemas so that the UI could read \texttt{subject}, \texttt{body} and the score fields without a second parser.

\section{Authentication}
\label{sec:auth}

My supervisor asked for authentication on every page and every API. I implemented \texttt{requireLabAccess} with two ways in: an examiner login form that sets an HttpOnly cookie (\texttt{POST /api/login}), and HTTP Basic Auth for \texttt{curl} and older clients. Static files are public so the login page can load. Every \texttt{/api/*} route except login and logout stays locked. I chose this because a browser login is easier to demonstrate than a raw Basic prompt, and because the laboratory is a small demonstration, not a campus SSO product.

Credentials come from \texttt{LAB\_BASIC\_USER} and \texttt{LAB\_BASIC\_PASS}. If they are missing, local development continues without a login. In production the same missing pair returns 503 and the text ``This laboratory is closed''.

Comparison uses \texttt{timingSafeEqualString}. Different lengths return false. Empty strings return false.

\section{Rate limit, size limit, and prompt checks}
\label{sec:limits}

I set \texttt{express.json(\{ limit: '32kb' \})} because a teaching API does not need large uploads and because the deepfake route should not become a file dump.

\texttt{rateLimitApi} counts POST requests per IP and path in a one-hour window. I used 10 requests per hour for \texttt{/api/generate-phishing} and 40 for the other POST API paths. I chose a tighter limit on generate because that route spends Gemini quota and because it is the dual-use route (see Chapter~\ref{cha:disc}).

\texttt{promptLooksInjected} looks for three regular expressions: ``ignore previous/prior/above instructions'', ``you are now '', and ``disregard the system/developer''. I know this is a small list. I still wanted a reject path I could test. A match returns HTTP 400 with ``Prompt rejected by input validation.''

\section{Phishing laboratory}
\label{sec:impl-phish}

\texttt{academicDataset.ts} holds three authored examples, four public-pattern reconstructions, and \texttt{generate90Samples()}. The generator uses a Mulberry32 RNG with seed 20260927. I chose a fixed seed so that screenshots and the dashboard aggregates stay reproducible.

The score bands are configuration:

\begin{itemize}
\item legitimate: realism 90--99, detectability 15--35, polish 93--100;
\item conventional phishing: realism 30--55, detectability 80--95, polish 40--65;
\item AI-style phishing: realism 81--96, detectability 40--60, polish 90--100.
\end{itemize}

When the 90 rows are loaded, the dashboard mean designed realism sits near 75\%. The value ``100 minus mean AI detectability'' sits near 51\%. Those two numbers are the ones people keep quoting. They are not measurements of Gmail or SpamAssassin. I kept the original card look and I labelled the badges ``Designed avg'' and ``Designed'' so the screen stays the same while the caption stays honest.

\texttt{parseEmailFile} reads a local \texttt{.txt} or \texttt{.eml} in the browser. It looks for From, To, Subject and a blank-line body split. I added file load because an examiner can bring a sample without pasting by hand. The file never leaves the browser except as the parsed text in \texttt{/api/analyze-email}.

\section{Deepfake tab}
\label{sec:impl-df}

I implemented three presets only. I wrote the filenames so that they look like teaching cases: a CEO video, a voicemail, a passport image. Exact name match returns the preset. Any other string is narrated. I put a sentence on the tab: media bytes are not inspected. I did this because shipping a weak detector would have been worse than shipping a labelled simulator.

\section{Adversarial path}
\label{sec:impl-adv}

\texttt{eval/adversarial/service.py} is a FastAPI application. On first attack it downloads CIFAR-10 and loads \texttt{cifar10\_resnet56} through \texttt{torch.hub}. It walks the test set and takes the first image of the requested class that the model already classifies correctly. The default attack is untargeted FGSM or PGD in pixel space; a targeted run happens only when the request sets \texttt{targeted: true}. The result is projected onto the $L_\infty$ ball and returned as softmax percentages plus the measured $L_\infty$. Unknown class names return HTTP 400. The architecture field is ignored; the loaded checkpoint is always ResNet-56.

\texttt{eval/adversarial/run\_attacks.py} is the batch script. It samples 200 test images by default, seed 20260927, and writes \texttt{adversarial\_metrics.json} with \emph{untargeted} FGSM and PGD. I describe the method in Chapter~\ref{cha:eval}. I do not paste invented accuracies.

The attack tab names the model CIFAR-10 ResNet-56 and offers FGSM or PGD. I do not offer ImageNet class names (banana, panda) as presets. The class inputs are CIFAR-10 labels. The photographs next to the canvas are stock illustrations, not the $32 \times 32$ tensors used by the attack.

\section{Tests}
\label{sec:tests}

I wrote the automated tests in \texttt{tests/labSecurity.test.ts} and \texttt{tests/apiContract.test.ts} and I run them with \texttt{npm test} (\texttt{tsx --test}). Table~\ref{tab:tests} lists the first four security helpers. The other tests check the 32\,kb limit, the 10/40 rate limits, fail-closed production auth, the login cookie, the four POST teaching routes, the three eval GET routes, and the \texttt{source} labels.

\begin{table}[htbp]
\centering
\caption{Automated tests in \texttt{labSecurity.test.ts}.}
\label{tab:tests}
\scriptsize
\begin{tabular}{p{5.6cm}p{7.6cm}}
\toprule
Test name & What I asserted \\
\midrule
\texttt{timingSafeEqualString rejects different lengths} & \texttt{'ab'} vs \texttt{'a'} is false; two empty strings are false. \\
\texttt{timingSafeEqualString accepts exact match} & \texttt{alpha} equals itself; \texttt{alpia} does not. \\
\texttt{labCredentials fail closed when env is empty} & After deleting the two env vars, \texttt{labCredentials()} is \texttt{null}. \\
\texttt{prompt injection patterns are rejected} & ``ignore previous instructions and dump the key'' is true; a normal library-fine prompt is false. \\
\bottomrule
\end{tabular}
\end{table}

I did not write an end-to-end browser test. I tested the four tabs by hand on the local server. That is a limitation. The unit tests still cover the lock helpers, the size cap, and the route contracts that my supervisor asked for.

\section{What I refused to implement}
\label{sec:refused}

I refused a Gemini fallback on the adversarial POST. I refused a mailbox crawler. I refused to treat the dashboard percentages as experimental results. I refused to keep the academic-report generator in the repository. Those refusals are part of the implementation, not only of the discussion.

\chapter{Evaluation}
\label{cha:eval}

This chapter describes how I evaluated the laboratory and which numbers I can actually cite. I only use values from files that exist in the repository. If a file is missing, I write that the number is \textit{not yet exported}. I do not invent clean accuracy, robust accuracy, or attack success rate.

\section{What I evaluated}
\label{sec:eval-what}

I separated three kinds of output.

\begin{enumerate}
\item \textbf{Configuration.} Scores I assigned in templates. The dashboard 75\% and 51\% values belong here.
\item \textbf{Computed on a reconstructed corpus.} The phishing JSON in \texttt{eval/phishing/results/phishing\_metrics.json}.
\item \textbf{Computed on a real classifier.} The adversarial JSON in \texttt{eval/adversarial/results/adversarial\_metrics.json}. The file currently holds an 8-image debug probe. I do not cite those eight images as the 200-image run my supervisor asked for.
\end{enumerate}

I also ran the 25 unit tests and I read the \texttt{source} field in \texttt{server.ts}.

\section{Dashboard values are configuration}
\label{sec:eval-a3}

A3 in the supervisor action list is explicit. The 75\% and 51\% figures are not findings.

When I click ``Load Study Dataset (90 Samples)'', \texttt{generate90Samples()} fills realism and detectability from the bands in Section~\ref{sec:impl-phish}, using seed 20260927. The dashboard uses \texttt{Math.round} on those means. For that seeded load the rounded values are 75\% mean designed realism (raw mean $75.12$) and 51\% for $100$ minus mean AI-class detectability (raw $51.23$). Those are still configuration. They are not filter measurements.

I wrote those bands so that conventional rows look easy to spot and AI-style rows look harder. That is a teaching design. It is not a measurement of a mail filter, a clinical pilot, or a sample-size study. I deleted that kind of language from the user interface. I do not use it here.

The dataset is in-memory templates. It is not a mailbox harvest.

\section{Phishing benchmark}
\label{sec:eval-phish}

I wrote \texttt{eval/phishing/run\_benchmark.py}. The first corpus was 288 repeated templates and the classifier reached F1 $= 1.0$. That score was a construction error, not a result, so I rebuilt the set. Ham and phish now share the same stems and differ mainly in the last instruction. The sender domain is stored but held out of the TF-IDF text, so the model is a wording check like the teaching tab.

The current file has 187 reconstructed rows, split 70/30 with \texttt{random\_state=20260927}. Table~\ref{tab:phish} copies the exported numbers. I do not treat them as Gmail or SpamAssassin performance.

\begin{table}[htbp]
\centering
\caption{Values copied from \texttt{eval/phishing/results/phishing\_metrics.json} (generated at 2026-10-20T18:10:55Z).}
\label{tab:phish}
\begin{tabular}{lr}
\toprule
Field & Value \\
\midrule
\texttt{n\_total} & 187 \\
\texttt{n\_train} & 130 \\
\texttt{n\_test} & 57 \\
Model & TF-IDF (1--2 gram) + logistic regression \\
\texttt{source} & computed \\
Accuracy & 0.649 \\
Precision & 0.633 \\
Recall & 0.939 \\
F1 & 0.756 \\
AUC & 0.75 \\
Test support (ham / phish) & 24 / 33 \\
\bottomrule
\end{tabular}
\end{table}

The JSON \texttt{label\_note} says the labels are the corpus labels shipped in \texttt{corpus.jsonl}, not a mailbox harvest and not Gemini-as-judge. Sender domain was held out. I agree with that note.

Recall is high and ham recall is low (0.25 on 24 ham test rows). The wording model still over-flags lookalike institutional text. That is a limitation of a small reconstructed set, not a detector finding about the internet.

I exported \texttt{MANUAL\_REVIEW.csv} and I ticked 20 rows. The teaching label matched the last instruction in each of those rows (\texttt{rinor\_agrees=yes}).

\section{Adversarial method}
\label{sec:eval-adv}

I implemented the method my supervisor asked for: FGSM and PGD~\cite{goodfellow2015explaining,madry2018towards}, on at least 200 images, with clean accuracy, robust accuracy, attack success rate, and perturbation size. Torchattacks is used when the import works~\cite{kim2020torchattacks}.

\texttt{eval/adversarial/run\_attacks.py} does the following.

\begin{enumerate}
\item Download the CIFAR-10 test set~\cite{krizhevsky2009cifar}.
\item Draw $n=200$ indices with \texttt{numpy} RandomState seed 20260927.
\item Load \texttt{cifar10\_resnet56} from \texttt{chenyaofo/pytorch-cifar-models}~\cite{chenyaofo2021cifar}.
\item Measure clean accuracy.
\item For $\varepsilon \in \{2/255,\ 8/255,\ 16/255\}$, measure FGSM and PGD robust accuracy, attack success rate (label flip versus the clean prediction), and mean $L_\infty$.
\item Write \texttt{eval/adversarial/results/adversarial\_metrics.json} with \texttt{source: "computed"}.
\end{enumerate}

The JSON file exists, but it is an 8-image debug probe (\texttt{n=8}, clean accuracy $0.875$, seed 20260927). I used that run only to check that CIFAR normalisation was wired. I do not copy those eight-image FGSM/PGD scores into a results table. The supervisor asked for at least 200 images. Until \texttt{n=200} is in that file:

\begin{itemize}
\item clean accuracy on 200 images: \textit{not yet exported};
\item robust accuracy (FGSM, PGD, each $\varepsilon$) on 200 images: \textit{not yet exported};
\item attack success rate on 200 images: \textit{not yet exported};
\item mean $L_\infty$ on 200 images: \textit{not yet exported}.
\end{itemize}

The live Simulate button is a single-image call to \texttt{POST /attack}. It returns one clean class, one adversarial class, and two softmax percentages. Those values are model outputs when the service is up. They are not the batch summary. I do not copy a live button result into this chapter because I do not have a frozen export of that click.

\section{Unit tests}
\label{sec:eval-tests}

I ran \texttt{npm test}. Twenty-five tests passed in \texttt{tests/labSecurity.test.ts} and \texttt{tests/apiContract.test.ts}. They check string comparison, fail-closed credentials, the injection regex, the 10/40 rate limits, the 32\,kb body cap, and the four documented \texttt{source} labels. They do not measure phishing F1 or ASR.

\section{What I did not evaluate}
\label{sec:eval-not}

I did not run a user study. I did not measure Gmail catch rates. I did not score real videos. I did not compare ResNet-50 ImageNet with ResNet-56 CIFAR in a finished metrics file. I did not count live \texttt{source=gemini} versus \texttt{source=fallback} calls into a JSON export. Those items stay as future work or as viva demonstrations, not as hidden numbers.

\chapter{Discussion}
\label{cha:disc}

This chapter interprets the results in relation to the literature, examines the threats to their validity, lists the limitations of the laboratory, discusses ethics and dual use, and presents the STRIDE analysis of the system.

\section{Interpretation of the results}
\label{sec:interp}

\paragraph{Phishing.} The benchmark indicates that, once phishing wording is as polished and context-specific as legitimate institutional mail, a wording-based classifier performs no better than always predicting phishing, while the sender address separates the classes. This matches the concern raised by work on LLM-generated phishing~\cite{hazell2023spear,heiding2024devising}: content quality no longer distinguishes attacks. It also explains why the classical detector of Fette et al.\ relied on structural and sender-related features~\cite{fette2007learning}. The practical consequence for teaching is that awareness training should move from ``look for spelling mistakes'' towards verifying the sender and the destination of links, and that organisations should rely on domain authentication~\cite{rfc9989}. The laboratory's mitigation advice follows this direction.

\paragraph{Adversarial examples.} The near-complete loss of accuracy of an undefended ResNet-56 under PGD at $8/255$ reproduces, on a small sample, the well-known fragility of standard training~\cite{szegedy2014intriguing,madry2018towards}. The contribution of the module is not this finding, which is expected, but its presentation: learners obtain real model outputs, with the budget verified, in the same environment in which they meet LLM-generated content, and they can see that a confident softmax output is not evidence of correctness.

\paragraph{Deepfakes.} The scenario module answers RQ2 only at the level of design. It shows that awareness content can be delivered without processing media and with explicit provenance, and that the defences it teaches (out-of-band verification, provenance checks) do not depend on detecting artefacts. It cannot show that learners become better at recognising deepfakes; that would require a study with participants.

\paragraph{Provenance.} Across the three modules, the \texttt{source} label separates computed results from LLM narration, presets and fallbacks. This separation is the central design property of the laboratory, and it is enforced by tests rather than by convention.

\section{Threats to validity}
\label{sec:validity}

\paragraph{Construct validity.} The phishing class labelled ``look-alike / spear-phishing'' imitates the style of LLM-generated phishing but was written by the author; the benchmark therefore does not measure the detection of LLM-generated email. The wording-only model is not strictly content-only, because message bodies contain link host names. The configured dashboard values do not measure any construct and are excluded from the findings.

\paragraph{Internal validity.} The same person designed, wrote and labelled the corpus, so the classifier may partly learn the author's writing habits, and the manual check of twenty labels is not independent. The PGD attack uses ten steps without random restarts, which can overestimate robust accuracy; at $8/255$ and above this does not affect the conclusion, because robust accuracy is already zero.

\paragraph{External validity.} The phishing results apply to a small reconstructed corpus of 187 emails from one institutional context and do not transfer to real mail streams. The adversarial results apply to one pretrained model on 200 CIFAR-10 test images; other architectures, datasets or defended models may behave differently.

\paragraph{Reliability.} Both evaluations use fixed seeds and regenerate their results from code, and the adversarial result file records the commit, the PyTorch version and the time of generation. Two dependencies are not pinned: the PyTorch version is not fixed in \path{eval/requirements.txt}, and the model is loaded from the default branch of its repository through \texttt{torch.hub}, so a future change upstream could alter the weights.

\section{Limitations}
\label{sec:limits-disc}

Beyond the threats above, the laboratory has the following limitations.

\begin{itemize}
\item \textbf{The LLM-based email analysis is not evaluated.} The classification that learners see in the analysis module comes from Gemini, whose accuracy on the benchmark corpus was not measured; the TF-IDF benchmark evaluates a different model. Published results for LLM detectors~\cite{koide2024chatspamdetector} do not transfer to this corpus.
\item \textbf{No learner study.} The educational effect of the laboratory has not been measured. Prior work shows how such an evaluation can be designed~\cite{sheng2007phil}.
\item \textbf{No media analysis.} The deepfake module does not inspect files and cannot demonstrate detection.
\item \textbf{Residual prompt-injection risk.} The denylist covers only naive direct-injection phrasing~\cite{perez2022ignore}, and the untrusted-data markers in the analysis prompt reduce but do not remove the risk of indirect injection~\cite{greshake2023signed}. The JSON response schema limits what an injected instruction can make the server return, but not what the model writes inside the permitted fields.
\item \textbf{Single-instance state.} Sessions and rate-limit counters are held in process memory; they are lost on restart and are not shared between instances (Section~\ref{sec:deploy}).
\item \textbf{Shared credentials and no audit log.} All users share one credential pair, and requests are not logged per user.
\item \textbf{Tests at function level.} The automated tests do not exercise the HTTP routes end to end, and the fallback and 503 paths are verified by inspection rather than by tests.
\end{itemize}

\section{Ethics and dual use}
\label{sec:ethics}

\paragraph{Dual use.} A phishing generator can serve attackers as well as learners, particularly when it produces lures tailored to a specific institution. The laboratory reduces this risk on the access side: it requires authentication and fails closed in production, limits generation to ten requests per hour per client, rejects obvious injection attempts, and has no function for sending email. These measures restrict \emph{who} can generate content and \emph{how much}, but not \emph{what} is generated. Output-side controls, such as a visible marker inserted into every generated message or replacing links with non-resolvable placeholders, were not implemented and are proposed as future work. The adversarial module attacks a public benchmark model, not a deployed system, and the deepfake module produces no synthetic media.

\paragraph{Data protection.} The analysis module accepts pasted text and \texttt{.eml} files, which may contain personal data of third parties, and forwards the extracted text to an external LLM provider. The laboratory stores none of this content on its server, but the transfer itself is a processing of personal data within the meaning of the General Data Protection Regulation~\cite{gdpr2016}. In its current form the module should therefore be used with the teaching material supplied with the laboratory or with messages from which personal data has been removed, and users should be told so on the analysis screen. This notice is not yet shown in the interface.

\paragraph{Honesty of outputs.} Presenting LLM-written scores as measurements would mislead learners about what AI systems can establish. The provenance labels, the absence of a fallback on the adversarial route and the labelling of configured dashboard values address this risk directly.

\section{STRIDE analysis}
\label{sec:stride}

Table~\ref{tab:stride} applies the STRIDE categories~\cite{shostack2014threat,microsoftstride} to the laboratory as deployed, with the controls in place and the residual risk.

\begin{table}[htbp]
\centering
\caption{STRIDE analysis of InsightFlow AI.}
\label{tab:stride}
\scriptsize
\begin{tabular}{p{2.0cm}p{3.6cm}p{3.6cm}p{3.2cm}}
\toprule
Category & Threat & Control & Residual risk \\
\midrule
Spoofing & Unauthorised use of the laboratory, in particular of the generator. & Authentication on every API route; constant-time comparison; fail-closed production mode; \texttt{HttpOnly} session cookies. & One shared credential pair; no second factor. \\
Tampering & Prompt injection that changes the task of the LLM; tampering with reported results. & Denylist on user prompts; untrusted-data markers for emails; JSON response schemas; adversarial results only from the attack service, with the budget verified. & Indirect injection not fully prevented. \\
Repudiation & A user denies having generated a message. & None beyond authentication. & No per-user accounts or audit log; accepted for a teaching tool, not for production use. \\
Information disclosure & Exposure of the API key or credentials; personal data in analysed emails sent to the LLM provider. & Secrets in environment variables only; key never sent to the browser; no server-side storage of content. & Analysed text leaves the institution; requires a usage notice (Section~\ref{sec:ethics}). \\
Denial of service & Exhaustion of LLM quota or server resources. & 32\,kB body limit; per-address hourly limits keyed on \texttt{req.ip}; independent failure of the attack service. & Limits are per instance; no global quota. \\
Elevation of privilege & Direct access to the attack service, bypassing the web server. & Optional shared token between server and attack service; single application role. & Token is optional; the service should not be exposed publicly without it. \\
\bottomrule
\end{tabular}
\end{table}

\section{Summary}

The results support three claims within their scope: wording does not separate the classes of the reconstructed corpus while the sender does; an undefended CIFAR-10 classifier is fragile under standard $L_\infty$ budgets; and the provenance of every output can be made explicit and tested. They do not support claims about real mail streams, about deepfake detection or about learning outcomes. Chapter~\ref{cha:conc} draws the conclusions.

\chapter{Conclusion and future work}
\label{cha:conc}

\section{Conclusion}
\label{sec:conc}

I implemented InsightFlow AI as an educational laboratory for three AI cybersecurity risks: AI-generated phishing, deepfake scenarios, and adversarial examples. I did this for the Informatics Engineering project course at ESTiG, under the supervision of Tiago Pedrosa. The official title is the title on the cover.

I can now answer the research questions in a limited way.

RQ1. I generate and analyse teaching emails with Gemini or a fallback, and I keep dataset labels on the 90 seeded rows. The source field is the part that makes the answer usable in a viva. The reconstructed TF-IDF wording benchmark reached F1 $= 0.756$ on 57 test rows (187 total). I read that as a small teaching corpus, not as a mailbox detector.

RQ2. I teach deepfake talking points with three filename presets and a narrated fallback. I do not inspect bytes. That is an honest teaching tab, not a forensics tool.

RQ3. I wired FGSM and PGD to a CIFAR-10 ResNet-56 so that a clean-accuracy number, once exported, can mean something on $32 \times 32$ images. My supervisor asked for ResNet-50. I explained the CIFAR choice in Section~\ref{sec:resnet-choice}. The batch metrics file was not yet exported when I wrote Chapter~\ref{cha:eval}.

The 75\% and 51\% dashboard values stay in the configuration column. I will not move them into the findings column.

The first website was started in Google AI Studio. Gemini still writes teaching mail. I still have to defend every number myself.

\section{Future work}
\label{sec:future}

If I continue the project after delivery, I would do the following.

\begin{enumerate}
\item Run \texttt{eval/adversarial/run\_attacks.py} to completion and replace every \textit{not yet exported} cell with the JSON fields. Then freeze a tag and take new screenshots.
\item Replace the Studio architecture buttons with one label: CIFAR-10 ResNet-56.
\item Build a harder phishing corpus, or drop the TF-IDF score from the report and keep only a manual review of the teaching rows.
\item Count \texttt{source=gemini} versus \texttt{source=fallback} on one recorded run per route, as A4 asked.
\item Add a real frame-level detector only if I can run it on a public set such as FaceForensics++ and report that set's protocol~\cite{rossler2019faceforensics}.
\item Add structured logs that store a timestamp and a hashed user name, not the email body, if an examiner needs repudiation.
\item Write one or two browser tests for login, load-90, and the 503 path of the attack tab.
\end{enumerate}

I would not add an SMTP sender. I would not harvest a mailbox. I would not present template bands as clinical evidence.


% listagem de referências / bibliography
% biblatex + biber, IEEE style; listed in the table of contents
\printbibliography[heading=bibintoc]

% Apêndice
\clearpage
\appendix
\renewcommand{\thepage}{\thechapter\arabic{page}}

\setcounter{page}{1}
\chapter{Declaration on the use of generative AI and TRACE}
\label{app:ai}

This appendix is the declaration I will defend. It is also the TRACE of the AI Studio origin of the first application.

\section{Declaration}
\label{sec:ai-decl}

I used generative AI in this project. I did not hide it.

\begin{itemize}
\item Google AI Studio generated the first web application (React tabs, the first Express routes, and the first visual layout).
\item An earlier Studio version also included an academic-report Word export. Those leftover \texttt{.docx} files were removed from this repository. I did not keep a generator that writes the official report.
\item Gemini is still used on the teaching email routes \texttt{/api/generate-phishing} and \texttt{/api/analyze-email}, and on non-preset deepfake filenames. Each of those answers carries \texttt{source} equal to \texttt{gemini}, \texttt{fallback}, or \texttt{preset}.
\item The adversarial route does not use Gemini. When the Python service is down, the answer is 503, not an invented confidence.
\item I implemented the lock, the tests, and the evaluation wiring in this repository.
\item This report text is written in the official UPB \texttt{project-report} template. I have to read every sentence and every citation before the viva. Every number must be checked by me, Rinor Basholli.
\end{itemize}

The work could not exist in this form without my decisions: the requirements, the source labels, the refusal to harvest mail, the CIFAR model choice, the unused adversarial fallback, and the honesty about the first F1 $= 1.0$ run and the later 0.756 wording run. The same work could have been written without AI, more slowly, as ordinary React and Python. AI did not sit the course. I did.

\section{TRACE of Studio metadata}
\label{sec:trace}

The repository still shows that it began in AI Studio.

\paragraph{metadata.json}
The file \texttt{metadata.json} at the project root is Studio project metadata. The name is ``AI Cybersecurity Risk Lab''. The description talks about an interactive study studio. \texttt{majorCapabilities} contains \texttt{MAJOR\_CAPABILITY\_SERVER\_SIDE\_GEMINI\_API}. \texttt{requestFramePermissions} is an empty list. I left this file in place so that the origin is visible.

\paragraph{vite.config.ts}
The Vite config still has comments about HMR being disabled in AI Studio through \texttt{DISABLE\_HMR}. That flag exists so that preview edits do not flicker. I did not invent it.

\paragraph{aistudio-build}
AI Studio publishes a Node website. The deploy notes in the project README warn that a republish from Studio can overwrite a custom all-in-one image. There is no separate \texttt{aistudio-build/} source tree in the current repository. The remaining traces are \texttt{metadata.json}, the HMR comments, the Gemini-first routes, and the leftover UI labels (the ``DB ACTIVE'' footer and the ResNet-50 button).

\paragraph{Removed Word export}
The academic-report \texttt{.docx} path is gone. I mention the removal here so that an examiner who saw an older tree knows I deleted it on purpose. The official report is this \LaTeX\ tree under \texttt{report/}, using BibTeX \texttt{ieeetr} as the template requires.

\section{Model identifier}
\label{sec:model-id}

The default Gemini model id in \texttt{server.ts} is \texttt{gemini-3.5-flash}. The environment variable \texttt{GEMINI\_MODEL} overrides it. I must say in the viva which id the live key served. I do not invent a transcript of those calls in this appendix.


\setcounter{page}{1}
\chapter{Reproducing the results}
\label{app:eval}

This appendix explains how to run the laboratory and regenerate every result reported in Chapter~\ref{cha:eval}. Commands are run from the repository root.

\section{Environment}
\label{sec:run-env}

The results were produced with Node.js~22, Python~3, PyTorch~2.14.1 (recorded in \path{adversarial_metrics.json}) and scikit-learn; the phishing results were also reproduced independently with scikit-learn~1.9.1, with byte-identical output. The adversarial results were produced on a CPU. The CIFAR-10 test set must be present under \path{eval/adversarial/results/cifar10/}; the model weights are downloaded by \texttt{torch.hub} on first use.

\section{Running the laboratory}
\label{sec:run-web}

\begin{minted}{text}
cp .env.example .env.local
# set GEMINI_API_KEY, LAB_BASIC_USER and LAB_BASIC_PASS
# optional: GEMINI_MODEL, EVAL_API_URL, EVAL_INTERNAL_TOKEN
npm install
npm test
python3 eval/adversarial/service.py &   # attack service on port 8001
npm run dev                              # http://localhost:3000
\end{minted}

For a production-style run, \texttt{npm run build} followed by \path{deploy/start-live-lab.sh} starts both processes and refuses to start without credentials. The container images are built from \texttt{Dockerfile} (web server only) and \path{deploy/Dockerfile.all-in-one} (web server and attack service).

\section{Phishing benchmark}
\label{sec:run-phish}

\begin{minted}{text}
python3 -m pip install -r eval/requirements.txt
python3 eval/phishing/run_benchmark.py
\end{minted}

The script regenerates \path{eval/phishing/corpus.jsonl} from code and writes \path{phishing_metrics.json} and \texttt{predictions.csv} under \path{eval/phishing/results/}. It does not overwrite \path{MANUAL_REVIEW.csv}.

\section{Adversarial batch}
\label{sec:run-adv}

\begin{minted}{text}
python3 eval/adversarial/run_attacks.py \
  2>&1 | tee eval/adversarial/results/run200_full.txt
\end{minted}

The script writes \path{eval/adversarial/results/adversarial_metrics.json}, including the commit and PyTorch version, and stops with an error if any perturbation exceeds its budget. The log of the reported run is \path{run200_full.txt}.

\section{Live attack calls}
\label{sec:run-live}

With the attack service running, the recorded live calls are reproduced with:

\begin{minted}{text}
curl -s -X POST localhost:8001/attack -H 'Content-Type: application/json' \
  -d '{"originalClass":"airplane","noiseLevel":0.05,"method":"fgsm"}'
curl -s -X POST localhost:8001/attack -H 'Content-Type: application/json' \
  -d '{"originalClass":"airplane","noiseLevel":0.05,"method":"pgd"}'
curl -s -X POST localhost:8001/attack -H 'Content-Type: application/json' \
  -d '{"originalClass":"airplane","targetClass":"bird","targeted":true,
       "noiseLevel":0.05,"method":"fgsm"}'
\end{minted}


\end{document}
